Benchmark scores for network intrusion detectors hide problems that appear only in deployment. We audit a decision tree on 54 NFStream flow features from labels to live blocking on an emulated Open vSwitch gateway, and condense the audit into four checks that others can run. \emph{Labels:} 91--98.5\% of flows in the CIC-IoT2023 web, brute-force and MITM captures are background traffic; training on file-name labels yields a true macro-F1 of 0.396 and flags 71.8\% of benign flows, against 0.725 and 3.8\% with evidence-based labels. In a robustness test on IoT-23 with the dataset authors' own Zeek labels (and our mapping of malware families to classes), the file-label tree flags 97--100\% of background flows in every malware capture. \emph{Splits:} on CIC-IoT2023 a random split hides that brute-force recall falls from 1.00 to 0.05 when the target changes from SSH to RTSP (three host pairs); on IoT-23 the temporal--random gap does not replicate; it is explained by missing classes and unstable under clustering. \emph{Live:} matching live to offline flows one by one exposed a forwarding pitfall, idle flow meters and an NFStream ordering defect that delayed verdicts by up to 82.5\,s. \emph{Response:} blocking on the first verdict blocked 60 of 1,674 eligible benign host-pair windows (3.6\%); requiring 10 verdicts blocked none but missed the Mirai and web attackers. Code, data and the audit scripts are public.
